% Generated draft: requires mathematical and source review.
\documentclass[11pt]{article}
\usepackage[margin=25mm]{geometry}
\usepackage{fontspec}
\setmainfont{Latin Modern Roman}
\usepackage{amsmath,amssymb,mathtools,unicode-math}
\setmathfont{Latin Modern Math}
\usepackage{enumitem,needspace,xurl,hyperref}
\setlength{\parindent}{0pt}
\setlength{\parskip}{5pt}
\setlength{\emergencystretch}{3em}
\newcommand{\N}{\mathbb{N}}
\newcommand{\Z}{\mathbb{Z}}
\newcommand{\Q}{\mathbb{Q}}
\newcommand{\R}{\mathbb{R}}
\newcommand{\C}{\mathbb{C}}
\newcommand{\F}{\mathbb{F}}
\newcommand{\A}{\mathbb{A}}
\newcommand{\PP}{\mathbb{P}}
\newcommand{\E}{\mathbb{E}}
\newcommand{\eps}{\varepsilon}
\newcommand{\dd}{\,\mathrm{d}}
\begin{document}
\section*{UnsolvedMath ID: 30005817}
\textbf{Source problem code:} catalogue entry
\subsection{Definitions and mathematical statement}
A four-dimensional polytope is a convex polytope \(P\subset\R^4\) with \(\dim P=4\). It is centrally symmetric if there exists \(a\in\R^4\) such that \(P=2a-P\). Write \(f_0(P)\) for its number of vertices and \(f_1(P)\) for its number of edges, meaning its one-dimensional faces. Central symmetry pairs the vertices, and a full-dimensional centrally symmetric polytope in \(\R^4\) has at least eight vertices. Accordingly, for every even integer \(n\geq8\), define
\[
c(n)=\max\{f_1(P):P\subset\R^4\text{ is a centrally symmetric four-dimensional convex polytope},\ f_0(P)=n\}.
\]
The principal problem is to determine \(c(n)\) for every such \(n\). The supplied original excerpt also asks whether
\[
c(n)=o(n^2),\qquad\text{that is,}\qquad
\lim_{\substack{n\to\infty\\n\text{ even}}}\frac{c(n)}{n^2}=0.
\]
This additional question is retained as supplied; it is absent from the cleaned statement, and its placement and intended notation require source review.

The original excerpt then introduces a separately numbered question, ``Problem 15 (Eran Nevo),'' about geometric triangles in \(\R^3\). Its stated quantity \(t(n)\) is the maximum number of triangles supported on \(n\) points such that any two distinct triangles intersect either in the empty set or in exactly one point that is a vertex of both triangles. Under the usual geometric interpretation, a triangle is the filled convex hull of three noncollinear points. That neighboring question asks whether \(t(n)=o(n^2)\). The excerpt does not explicitly settle whether every supporting point must be used or explicitly define nondegeneracy; these conventions require confirmation if this separate question is formalized.

The appended remarks state that, for combinatorial or piecewise-linear triangles, ``\(n^2\) triangles are possible,'' using Steiner triple systems embedded piecewise linearly in \(\R^3\), and attribute the lower bound \(t(n)=\Omega(n^{3/2})\) to K\'arolyi and Solymosi. The first remark is preserved as source wording: it does not specify whether \(n^2\) denotes an exact count or quadratic order. These neighboring statements are supplied context, not assertions established by this draft. They are not incorporated into the definition of \(c(n)\), and the attribution to Eran Nevo is attached only to the explicitly numbered triangle problem.
\subsection{Short English statement}
For each even \(n\geq8\), determine the largest possible number of edges of a centrally symmetric convex four-dimensional polytope with exactly \(n\) vertices. The raw excerpt additionally asks whether this maximum is \(o(n^2)\), then includes a separately numbered problem about triangles in \(\R^3\); the boundary between these source passages needs review. The supplied open-status assessment is dated curation, not an independently verified conclusion.
\subsection{Sources}
\begin{itemize}
\item Supplied local UnsolvedMath registry record: permanent ID \(30005817\), title ``Edges of Centrally Symmetric Four-Polytopes,'' source problem code \texttt{catalogue entry}, in the Oberwolfach Reports Open Problems collection. The supplied record labels the problem open and the statement \texttt{corrected\_verified}; its verification note says that the statement was reconstructed and checked against the official report. These are supplied registry assertions, not verification performed for this draft. No proposer is supplied for the polytope problem; the record gives the year 2024.
\item \emph{Discrete Geometry} (2024), Oberwolfach Reports, DOI 10.4171/owr/2024/3, pp.~137--202. Supplied report provenance: \url{https://doi.org/10.4171/owr/2024/3}.
\item Novik, \emph{A tale of centrally symmetric polytopes and spheres}. Supplied literature source: \url{https://sites.math.washington.edu/~novik/publications/cs-survey.pdf}.
\item \emph{Oberwolfach Report 3/2024}. Supplied literature source: \url{https://ems.press/content/serial-article-files/48651}.
\item The supplied literature assessment was checked on \texttt{2026-08-22}. It labels the problem open and states that the extremal edge count remains unknown asymptotically and that current surveys explicitly describe even its asymptotic behavior as unknown. This is supplied dated curation derived from the Oberwolfach material, not an independently verified literature assessment; its time-sensitive status and cited claims require independent checking.
\item The appended triangle passage explicitly attributes its Problem 15 to Eran Nevo and cites K\'arolyi and Solymosi, \emph{Almost Disjoint Triangles in 3-Space}, Discrete \& Computational Geometry 28 (2002), pp.~577--583, for the stated lower bound. No URL for this article is supplied. This attribution belongs to the neighboring triangle passage and does not identify the proposer of the polytope problem.
\end{itemize}
\end{document}
